\documentclass[11pt]{article}
\usepackage[margin=0.78in]{geometry}
\usepackage{amsmath,amssymb,bm}
\usepackage{booktabs}
\usepackage{enumitem}
\usepackage{hyperref}
\hypersetup{colorlinks=true, linkcolor=black, urlcolor=black}
\setlength{\parindent}{0pt}
\setlength{\parskip}{0.55em}

\title{Cyrus Guided Notes: Stanford CS231n Deep Learning for Computer Vision}
\author{Stanford course pack}
\date{2026-05-15}

\begin{document}
\maketitle

\section*{Course source}
\textbf{Official source:} \url{https://cs231n.stanford.edu/}

\textbf{Why this course is in the system.} CNNs, optimization, backpropagation, visual representation, and spatial-intelligence transfer.

\textbf{Learning output.} Already connected to guided PDF, notes index, and assignment links.

\section*{Prerequisite checkpoint}
\begin{itemize}[leftmargin=1.4em]
  \item Linear algebra
  \item Calculus
  \item Python and NumPy
\end{itemize}

\section*{Formula checkpoint}
Do not memorize these first. Read each formula as a sentence: what changes, what is observed, what is optimized, and what output it produces.
\begin{align*}
  p_k=\frac{e^{s_k}}{\sum_j e^{s_j}}\\
  L=-\log p_y\\
  (I*K)_{i,j}=\sum_{m,n}I_{i+m,j+n}K_{m,n}\\
  \frac{\partial L}{\partial W}=\frac{\partial L}{\partial s}\frac{\partial s}{\partial W}
\end{align*}

\section*{Visual page}
Draw pixels -> convolution -> activation -> class score -> loss.

\section*{GoodNotes pages}
\begin{itemize}[leftmargin=1.4em]
  \item Page ST-DL001: softmax and cross entropy
  \item Page ST-DL002: convolution
  \item Page ST-DL003: backprop
\end{itemize}

\section*{Web interaction}
A small convolution kernel visual that updates feature responses.

\section*{Obsidian and Notion}
\begin{tabular}{p{0.22\linewidth}p{0.68\linewidth}}
\toprule
Layer & Output \\
\midrule
Obsidian & Stanford Spatial AI -> CS231n \\
Notion & Assignment and formula-progress status. \\
\bottomrule
\end{tabular}

\section*{First study move}
Open the official source, copy the prerequisite checkpoint into GoodNotes, then write one formula in your own words before watching or reading more.

\end{document}
